\clearpage
\section{Appendix: Scope and Qualifications}
The target is the U.S. Twenty-seventh Sunday in Ordinary Time, Year A, 4 October 2026: green, Lectionary 139, permanent formula PC-S53-A. The governing books are the \work{Roman Missal, Third Edition}, U.S. 2011, based on the 2008 emended Latin edition, and the U.S. \work{Lectionary for Mass}, second typical edition (1998/2001), Volume I. The national calendar assigns this Sunday rather than the occurring Francis of Assisi memorial; local dedication, patronal or religious-calendar exceptions are not resolved.

Both Missal Communion antiphons are studied; neither is a local selection. Permitted entrance, psalm and Communion chant routes remain open, as do preparation music and coupled Preface/Eucharistic Prayer choices. No unique Week 27 Preface, Offertory antiphon, Sequence, Eucharistic Prayer insertion or distinctive conclusion is appointed in the inspected formulary. GIRM 63(c)'s one-reading condition does not authorize omission of the acclamation at this three-reading Sunday. Sprinkling and optional concluding blessings remain unselected. Ordinary Sunday requirements are not additional closed proper texts.

The concise argument derives from the same owner's expansive study and its two developed interpretations. Cross-element readings are editorial syntheses; no Father is credited with this complete modern formulary. Isaiah and Matthew are correlated selections; Philippians is semi-continuous and the Missal prayers recur across cycles. The historical English witnesses of Augustine's psalm and O'Sullivan's Bellarmine are abridged. No fresh critical-language collation is claimed. Direct, securely attributed saintly exposition of the precise Esther and Lamentations loci was not established; their biblical contexts support the discussion. Aquinas and Leo illuminate sacramental meaning rather than expound the modern prayers.

Scriptural wording where discussed follows public-domain Douay--Rheims Challoner study English, \textbf{not the approved English proclaimed at Mass}; the U.S. Lectionary controls that text. Protected Lectionary, ICEL and full Latin prayer bodies are not reproduced or reconstructed. The Week 27 evidence retains its limits: a 2002 digital Latin witness, official Antiphonary and FDLC excerpts, and a bounded 2008 variation notice do not constitute a complete 2008 or U.S. 2011 altar-book collation. The exact offerings English remains uncollated; the FDLC Collect's older conclusion is not adopted. No new prayer translation is supplied.

The chronological proposals distinguish composition, traditional attribution and event; the Catholic Encyclopedia witnesses are historical judgments, not current consensus. The source study and research records preserve the detailed search bounds, source bindings and unresolved evidence. This work is a study aid, not an official liturgical book or ecclesiastical approval of the interpretations.

\section{References}
\begin{itemize}
\item \work{Holy Bible}, Douay--Rheims, Challoner revision, Project Gutenberg 1581: Esther 4,13--14; Isaiah 5; Psalm 79; Lamentations 3; Matthew 21; John 15; 1 Corinthians 10; Philippians 4. Public-domain study witness.
\item Augustine, \work{Tractates on John} 86--87, trans. John Gibb and James Innes, NPNF1-7; \work{Exposition of Psalm} 79, §§6--11, NPNF1-8 (English heading Psalm 80). CCEL historical-English witnesses; psalm abridged.
\item John Chrysostom, \work{Homilies on Matthew} 68.1--2, NPNF1-10; \work{Homilies on Philippians} 14, NPNF1-13; \work{Homilies on First Corinthians} 24.4, NPNF1-12. CCEL historical-English witnesses.
\item Robert Bellarmine, \work{A Commentary on the Book of Psalms}, trans. John O'Sullivan (1866), Psalm 79. Abridged English, e-Catholic 2000 witness.
\item Thomas Aquinas, \work{Summa theologiae} III, q. 79, a. 1, body and replies, Dominican English, Gutenberg 19950; Leo the Great, Sermon 63, VII and note 1063, trans. Charles Lett Feltoe, NPNF2-12.
\item \work{General Introduction to the Lectionary}, 23,66--68,106--107; U.S. \work{General Instruction of the Roman Missal}, 48,51,53--54,61--71,74,87--90,364--365; USCCB \work{Liturgical Calendar 2026}, pp. 5,40, and \href{https://bible.usccb.org/bible/readings/100426.cfm}{4 October readings}, Lectionary 139. Dated witnesses checked in the source research, 28 September 2026.
\item \work{Missale Romanum} (2002), digital reproduction, physical pp. 290--291; ICEL \work{Antiphonary} (2010), p. 80; FDLC \work{Mystagogical Reflections}, Ordinary Time 24--30 (2016), pp. 12--13; \work{Notitiae} 44 (2008), pp. 367--372, especially 371. Their bounded roles are stated above.
\item \work{Catholic Encyclopedia}: McMahon, \work{Esther} (1909); Souvay, \work{Isaias} (1910); Prat, \work{St. Paul} (1911); Vander Heeren, \work{Epistle to the Philippians} (1911); Aherne, \work{Epistles to the Corinthians} (1908); Fonck, \work{Gospel of St. John} (1910); Maas, \work{Chronology of the Life of Jesus Christ} (1910); Durand, \work{The Synoptics} (1912); Jacquier, \work{Gospel of St. Matthew} (1911). NABRE introductions to Psalms and Lamentations supply the stated critical fallback horizons; protected Bible wording is not reproduced.
\end{itemize}
